\documentclass[11pt,a4paper]{article}

\usepackage[margin=25mm]{geometry}
\usepackage[T1]{fontenc}
\usepackage[utf8]{inputenc}
\usepackage{lmodern}
\usepackage{microtype}
\usepackage{amsmath,amssymb}
\usepackage{booktabs}
\usepackage{graphicx}
\usepackage{float}
\usepackage{siunitx}
\usepackage{xcolor}
\usepackage{hyperref}
\usepackage{enumitem}
\usepackage{listings}
\usepackage{longtable}

\hypersetup{
  colorlinks=true,
  linkcolor=blue!55!black,
  citecolor=blue!55!black,
  urlcolor=blue!55!black,
  pdftitle={Phase 1 Verification of a PyTorch pICNN Hyperelastic Material Model Coupled to NGSolve},
  pdfauthor={Vignesh Ravi}
}

\sisetup{
  scientific-notation=true,
  round-mode=places,
  round-precision=4
}

\lstset{
  basicstyle=\ttfamily\small,
  frame=single,
  breaklines=true,
  columns=fullflexible,
  showstringspaces=false
}

\newcommand{\F}{\mathbf{F}}
\newcommand{\C}{\mathbf{C}}
\newcommand{\Pstress}{\mathbf{P}}
\newcommand{\Aten}{\mathbb{A}}
\newcommand{\Iten}{\mathbf{I}}
\newcommand{\R}{\mathbf{R}}
\newcommand{\placeholder}[1]{\textcolor{red!70!black}{\textbf{[#1]}}}

\newcommand{\optionalfigure}[3]{%
  \begin{figure}[H]
    \centering
    \IfFileExists{#1}{%
      \includegraphics[width=0.82\textwidth]{#1}
    }{%
      \fbox{\parbox[c][45mm][c]{0.78\textwidth}{%
        \centering Optional figure file not found:\\[1mm]
        \texttt{\detokenize{#1}}}}
    }
    \caption{#2}
    \label{#3}
  \end{figure}
}

\title{%
  \textbf{Phase 1 Verification of a PyTorch pICNN\\
  Hyperelastic Material Model Coupled to NGSolve}
}
\author{%
  Vignesh Ravi\\
  Computational Materials Science\\[1mm]
  Supervisor: Stefan Prüger
}
\date{\today}

\begin{document}

\maketitle

\begin{abstract}
This report documents Phase~1 of a proof-of-concept framework for coupling a
PyTorch-based polyconvex input-convex neural network (pICNN) hyperelastic
material model to NGSolve. The surrogate was trained against a compressible
neo-Hookean reference material and evaluated in a total-Lagrange
finite-element formulation. First-Piola stress and the full material tangent
were computed using PyTorch automatic differentiation, without finite
differences in the constitutive update. Material values were transferred to
NGSolve using \texttt{IntegrationRuleSpace}.

Verification covered the model interface, stress-free reference
normalization, agreement with the analytical neo-Hookean response, tangent
consistency, nonlinear element convergence, a multi-element affine patch
test, and matched analytical-versus-neural finite-element problems. All
predefined Phase~1 acceptance criteria were satisfied. In the matched
finite-element benchmark, the free-displacement, deformation-gradient, and
reaction-vector relative errors were \SI{1.376}{\percent},
\SI{0.3504}{\percent}, and \SI{13.60}{\percent}, respectively. These results
validate the PyTorch--NGSolve bridge within the stated Phase~1 scope. They do
not yet validate the particle RVE formulation planned for Phase~2.
\end{abstract}

\tableofcontents
\clearpage

\section{Objectives and scope}

\subsection{Phase 1 objectives}

The objectives of Phase~1 were to:
\begin{enumerate}[label=\arabic*.]
  \item implement a pICNN strain-energy model with non-negative constrained
        weights;
  \item use the feature vector
        $\left[I_1,I_2,J,-\ln J\right]$ on the admissible domain $J>0$;
  \item enforce zero energy and zero stress at the undeformed reference state;
  \item compute first-Piola stress and the full material tangent using
        automatic differentiation;
  \item transfer the constitutive response through NGSolve
        \texttt{IntegrationRuleSpace};
  \item demonstrate consistent Newton convergence and finite-element
        equilibrium; and
  \item compare neural constitutive and finite-element responses with the
        analytical neo-Hookean reference.
\end{enumerate}

\subsection{Items outside the Phase 1 scope}

Phase~1 does not validate particle geometry, particle--matrix interfaces,
periodic RVE constraints, Hill--Mandel homogenization, rigid-particle
kinematics, a mixed near-incompressible formulation, volume-fraction context
inputs, or extrapolation outside the Phase~1 training domain. These items are
reserved for Phase~2 and require separate verification.

\section{Analytical reference material}

The Phase~1 target is the compressible neo-Hookean strain-energy density
\begin{equation}
  \Psi_{\mathrm{NH}}(\F)
  =
  \frac{\mu}{2}(I_1-3)
  -\mu\ln J
  +\frac{\lambda}{2}(\ln J)^2,
  \label{eq:nh-energy}
\end{equation}
where
\begin{equation}
  \C=\F^\mathsf{T}\F,\qquad
  I_1=\operatorname{tr}(\C),\qquad
  J=\det\F>0.
\end{equation}
The analytical first-Piola stress is
\begin{equation}
  \Pstress_{\mathrm{NH}}
  =
  \mu\left(\F-\F^{-\mathsf{T}}\right)
  +\lambda\ln(J)\F^{-\mathsf{T}}.
  \label{eq:nh-stress}
\end{equation}
The nondimensional Phase~1 parameters are
\begin{equation}
  \mu=1,\qquad \lambda=10.
\end{equation}
These values define a well-conditioned proof-of-concept material and are not
intended as a calibrated PDMS model.

\section{pICNN formulation}

\subsection{Invariant features}

The right Cauchy--Green tensor and its first two principal invariants are
\begin{align}
  \C &= \F^\mathsf{T}\F,\\
  I_1 &= \operatorname{tr}(\C),\\
  I_2 &=
  \frac{1}{2}\left[
    I_1^2-\operatorname{tr}(\C^2)
  \right].
\end{align}
The neural feature vector is
\begin{equation}
  \mathbf{y}(\F)
  =
  \left[I_1,I_2,J,-\ln J\right].
  \label{eq:features}
\end{equation}

The ICNN uses softplus activations and positive effective weights. For a
stored unconstrained parameter $\theta$, the corresponding effective weight
is
\begin{equation}
  W^+ = \operatorname{softplus}(\theta)>0.
\end{equation}
Feature-to-hidden, hidden-to-hidden, and output weights are constrained in
this manner. The architecture is constructed to preserve convexity and
monotonicity with respect to the selected polyconvex features, supporting a
polyconvex energy representation on $J>0$. Polyconvexity is not interpreted
here as a guarantee that the tangent is positive definite for every arbitrary
deformation.

\subsection{Stress-free reference normalization}

Let $\Psi_{\mathrm{raw}}(\F)$ denote the raw network output. The normalized
energy is
\begin{equation}
  \widehat{\Psi}(\F)
  =
  \Psi_{\mathrm{raw}}(\F)
  -\Psi_{\mathrm{raw}}(\Iten)
  -\Pstress_{\mathrm{raw}}(\Iten):(\F-\Iten),
  \label{eq:normalization}
\end{equation}
where
\begin{equation}
  \Pstress_{\mathrm{raw}}(\Iten)
  =
  \left.
  \frac{\partial\Psi_{\mathrm{raw}}}{\partial\F}
  \right|_{\F=\Iten}.
\end{equation}
Equation~\eqref{eq:normalization} gives
\begin{equation}
  \widehat{\Psi}(\Iten)=0,
  \qquad
  \left.
  \frac{\partial\widehat{\Psi}}{\partial\F}
  \right|_{\F=\Iten}
  =\mathbf{0}.
\end{equation}
The subtracted term is affine in $\F$ and therefore does not change the
second derivative of the raw energy.

\section{Automatic differentiation and tangent convention}

The first-Piola stress and material tangent are
\begin{align}
  P_{iJ}
  &=
  \frac{\partial\widehat{\Psi}}{\partial F_{iJ}},\\
  A_{iJkL}
  &=
  \frac{\partial P_{iJ}}{\partial F_{kL}}
  =
  \frac{\partial^2\widehat{\Psi}}
       {\partial F_{iJ}\partial F_{kL}}.
  \label{eq:tangent}
\end{align}
Because $\F$ is unsymmetric, all nine components are retained. The row-major
matrix representation used during assembly is
\begin{equation}
  A^{9\times9}_{\,3i+J,\;3k+L}=A_{iJkL}.
  \label{eq:flattening}
\end{equation}
A $6\times6$ minor-symmetric Voigt tangent would be appropriate for a
formulation based directly on symmetric strain and stress measures. It is not
the correct direct representation for the derivative
$\partial\Pstress/\partial\F$ used here.

No finite differences are used to calculate stress or tangent during
training or finite-element assembly. Central differences are used only in
Verification~V3 as an independent check of the autograd tangent.

\section{PyTorch--NGSolve bridge}

The constitutive bridge performs the following operations during a Newton
iteration:
\begin{enumerate}[label=\arabic*.]
  \item construct $\F=\Iten+\nabla\mathbf{u}$;
  \item interpolate its nine components into an
        \texttt{IntegrationRuleSpace};
  \item read each component separately and stack the values into an
        $N\times3\times3$ PyTorch tensor;
  \item evaluate $\Pstress$ and $\Aten$ with automatic differentiation;
  \item write the nine stress and 81 tangent components back
        component-by-component; and
  \item assemble the residual and tangent forms.
\end{enumerate}

The total-Lagrange residual and tangent have the forms
\begin{align}
  R_a &=
  \int_{\Omega_0}
  \Pstress:\nabla N_a\,\mathrm{d}V,\\
  K_{ab} &=
  \int_{\Omega_0}
  \nabla N_a:\Aten:\nabla N_b\,\mathrm{d}V.
\end{align}
The Newton residual norm is evaluated on free degrees of freedom; constrained
entries represent reaction forces rather than equilibrium equations.

NGSolve compound \texttt{IntegrationRuleSpace} storage is component-major.
The production bridge avoids layout ambiguity by reading and writing
individual components. Directly reshaping a complete compound vector as
point-major is not valid.

\section{Training procedure}

The model was trained using an energy--stress Sobolev objective,
\begin{equation}
  \mathcal{L}
  =
  \lambda_\Psi
  \operatorname{MSE}
  \left(\widehat{\Psi},\Psi_{\mathrm{NH}}\right)
  +
  \lambda_P
  \operatorname{MSE}
  \left(\Pstress,\Pstress_{\mathrm{NH}}\right),
  \label{eq:loss}
\end{equation}
with
\begin{equation}
  \lambda_\Psi=1,\qquad \lambda_P=10.
\end{equation}
Training used CPU PyTorch and float64 constitutive calculations. Adam was
used with an initial learning rate of $2\times10^{-3}$, followed by checkpoint
refinement at $1\times10^{-3}$.

\subsection{Initialization diagnosis and correction}

The original unconstrained weight parameters were initialized near zero.
Since the effective constrained weights are generated by softplus, this gave
\begin{equation}
  \operatorname{softplus}(0)\approx0.693,
\end{equation}
which is not a small initial weight. Across multiple 32-unit layers, this
produced an excessively stiff network and a very large initial loss.

\begin{table}[H]
  \centering
  \caption{Training behavior before and after correcting positive-weight
  initialization.}
  \label{tab:init}
  \begin{tabular}{lS}
    \toprule
    Quantity & {Value}\\
    \midrule
    Original initial loss & 1.173628e9\\
    Original final loss after 1500 epochs & 1.128706e7\\
    Corrected initial loss & 4.456922e1\\
    Corrected mean loss over epochs 1400--1499 & 2.772083\\
    \bottomrule
  \end{tabular}
\end{table}

The corrected effective initialization values were approximately 0.05 for
feature weights, 0.02 for hidden-to-hidden weights, and 0.05 for output
weights. The failed checkpoint was retained separately and was not used for
verification.

\optionalfigure
  {figures/training_loss.pdf}
  {Training and refinement loss history on a logarithmic scale.}
  {fig:training-loss}

\section{Verification methodology and results}

\subsection{Acceptance matrix}

\begin{table}[H]
  \centering
  \small
  \caption{Phase 1 verification matrix.}
  \label{tab:verification}
  \begin{tabular}{p{0.08\textwidth}p{0.38\textwidth}p{0.35\textwidth}c}
    \toprule
    Test & Purpose & Acceptance criterion & Result\\
    \midrule
    V1 & Interface, tensor shapes, dtype and derivative finiteness
       & Correct shapes and finite float64 outputs & PASS\\
    V2 & Stress-free reference state
       & Energy and stress below scale-aware tolerance & PASS\\
    V2b & Analytical constitutive response
       & Energy $<10\%$, stress $<10\%$, tangent $<25\%$ & PASS\\
    V3 & Autograd tangent verification
       & Best relative central-difference error $<10^{-6}$ & PASS\\
    V4 & Neural nonlinear element solve
       & Final free residual $<10^{-8}$ & PASS\\
    V5 & Multi-element affine patch
       & Gradient, stress variation and residual $<10^{-10}$ & PASS\\
    V6 & Matched analytical and neural FE problems
       & Field and reaction errors $<15\%$ & PASS\\
    \bottomrule
  \end{tabular}
\end{table}

\subsection{V1: model interface}

The verified public material interface accepts deformation gradients through
\texttt{model.energy(F)}. Precomputed invariant features are passed to the
wrapped ICNN rather than to the normalized constitutive interface. Energy,
stress and tangent tensors had the expected shapes and finite float64 values.
The stress-only training path agreed with the stress returned by the full
stress-and-tangent path.

\subsection{V2: reference state}

\begin{table}[H]
  \centering
  \caption{Reference-state verification.}
  \label{tab:v2}
  \begin{tabular}{lS}
    \toprule
    Metric & {Value}\\
    \midrule
    $\left|\widehat{\Psi}(\Iten)\right|$ & 0.0\\
    $\max\left|P_{iJ}(\Iten)\right|$ & 7.276e-12\\
    Raw reference stress scale & 4.866e4\\
    Scale-aware stress tolerance & 1.081e-10\\
    \bottomrule
  \end{tabular}
\end{table}
The normalized stress residual is approximately float64 cancellation error
relative to the raw stress scale and does not indicate physical prestress.

\subsection{V2b: analytical constitutive comparison}

The trained checkpoint was evaluated on deterministic deformation gradients
not used for training. The resulting relative errors were:
\begin{table}[H]
  \centering
  \caption{Neural versus analytical neo-Hookean constitutive response. Replace
  the red placeholders with the exact values from
  \texttt{reports/logs/v2b\_neohooke.txt}.}
  \label{tab:v2b}
  \begin{tabular}{lcc}
    \toprule
    Quantity & Relative error & Acceptance threshold\\
    \midrule
    Energy & \placeholder{insert V2b energy error} & 0.10\\
    First-Piola stress & \placeholder{insert V2b stress error} & 0.10\\
    Full tangent & \placeholder{insert V2b tangent error} & 0.25\\
    \bottomrule
  \end{tabular}
\end{table}

\subsection{V3: tangent verification}

The autograd tangent was compared with a central-difference derivative of
the first-Piola stress over multiple perturbation sizes. The expected
V-shaped error curve results from the competition between finite-difference
truncation error and floating-point roundoff. The best measured relative
error was \placeholder{insert best V3 relative error}, satisfying the
$10^{-6}$ criterion.

\optionalfigure
  {figures/tangent_fd_check.pdf}
  {Relative central-difference tangent error versus perturbation size.}
  {fig:tangent-check}

\subsection{V4: nonlinear neural element solve}

The trained neural material was inserted into the NGSolve Newton solver using
the exact autograd tangent. The final free residual was
\placeholder{insert V4 final residual}. The residual history was
\placeholder{insert compact V4 residual sequence}. The solve satisfied the
$10^{-8}$ convergence criterion.

\optionalfigure
  {figures/v4_newton_residual.pdf}
  {V4 neural finite-element Newton residual history.}
  {fig:v4-newton}

\subsection{V5: multi-element affine patch test}

A mixed stretch/shear affine displacement field was prescribed over a mesh
with free interior degrees of freedom. A first-order conforming
discretization should reproduce the constant deformation gradient exactly.

\begin{table}[H]
  \centering
  \caption{V5 affine patch-test results.}
  \label{tab:v5}
  \begin{tabular}{lS}
    \toprule
    Metric & {Value}\\
    \midrule
    Mesh elements & 206\\
    Material points & 2884\\
    Free degrees of freedom & 30\\
    $\max|\F-\F_{\mathrm{target}}|$ & 4.441e-16\\
    Relative stress variation & 4.774e-15\\
    Free residual norm & 2.083e-16\\
    \bottomrule
  \end{tabular}
\end{table}
The machine-precision results verify affine completeness, integration-point
data transfer, tensor ordering, constitutive uniformity, and free-DOF
equilibrium.

\subsection{V6: matched analytical and neural FE problems}

The analytical and neural materials were solved on the same 206-element mesh
under identical boundary conditions and loading. Both problems converged in
five Newton iterations.

\begin{table}[H]
  \centering
  \caption{V6 nonlinear convergence.}
  \label{tab:v6-convergence}
  \begin{tabular}{lcc}
    \toprule
    Quantity & Analytical material & Neural material\\
    \midrule
    Newton iterations & 5 & 5\\
    Final free residual & $3.657\times10^{-14}$ & $1.128\times10^{-14}$\\
    \bottomrule
  \end{tabular}
\end{table}

\begin{table}[H]
  \centering
  \caption{V6 neural-versus-analytical FE comparison.}
  \label{tab:v6-errors}
  \begin{tabular}{lcc}
    \toprule
    Comparison & Relative error & Acceptance threshold\\
    \midrule
    Free displacement & 0.01376 & 0.15\\
    Deformation gradient & 0.003504 & 0.15\\
    Constrained reaction vector & 0.1360 & 0.15\\
    \bottomrule
  \end{tabular}
\end{table}

The deformation field is reproduced closely. The reaction vector is more
sensitive to constitutive approximation and gives the largest observed
discrepancy. It satisfies the predefined Phase~1 threshold but remains an
important target for future surrogate refinement.

\optionalfigure
  {figures/v6_newton_comparison.pdf}
  {Analytical and neural V6 Newton residual histories.}
  {fig:v6-newton}

\section{Limitations}

\begin{enumerate}[label=\arabic*.]
  \item The training and validation domain is restricted to moderate
        deformations around the reference configuration.
  \item Phase~1 uses a homogeneous synthetic neo-Hookean target rather than
        experimental or RVE-generated data.
  \item The reaction-vector discrepancy of \SI{13.60}{\percent} is larger than
        the displacement and deformation-gradient discrepancies.
  \item The nondimensional material parameters are not calibrated to PDMS.
  \item Particles, interfaces, periodicity and volume-fraction dependence have
        not yet been introduced.
  \item Near-incompressibility has not yet been treated using a mixed
        displacement--pressure formulation.
  \item Central differences are used only as an independent verification
        method and are not part of the constitutive solver.
  \item Passing the Phase~1 tests does not establish accuracy under arbitrary
        extrapolation.
\end{enumerate}

\section{Conclusion and Phase 2 gate}

Phase~1 demonstrates an end-to-end coupling between a PyTorch pICNN
hyperelastic model and NGSolve. The implementation provides a stress-free
normalized energy, first-Piola stress, and a full $9\times9$ total-Lagrange
tangent through automatic differentiation. Integration-point transfer was
verified through a multi-element affine patch test, and the consistent
tangent produced converged Newton iterations in both analytical and neural
finite-element problems.

The neural model satisfied the predefined constitutive and finite-element
acceptance criteria. In the matched FE benchmark, the free-displacement and
deformation-gradient errors were \SI{1.376}{\percent} and
\SI{0.3504}{\percent}, while the reaction-vector error was
\SI{13.60}{\percent}. The reaction result is acceptable for the current
proof-of-concept threshold but identifies a target for further refinement.

On this basis, the PyTorch--NGSolve bridge is considered validated for the
scope of Phase~1. Phase~2 may proceed to periodic RVE generation for rigid
particles embedded in a soft matrix, subject to separate verification of
periodic constraints, rigid-particle kinematics, mixed near-incompressible
behavior, Hill--Mandel consistency, and the homogenized tangent.

\appendix

\section{Reproducibility information}

\begin{table}[H]
  \centering
  \caption{Software environment. Complete this table from
  \texttt{reports/logs/environment.txt}.}
  \begin{tabular}{ll}
    \toprule
    Item & Version\\
    \midrule
    Python & 3.12 (Anaconda environment)\\
    PyTorch & \placeholder{insert exact version}\\
    NGSolve & 6.2.2607\\
    NumPy & \placeholder{insert exact version}\\
    Execution device & CPU\\
    Constitutive dtype & float64\\
    \bottomrule
  \end{tabular}
\end{table}

The validated model checkpoint is
\begin{center}
  \texttt{checkpoints/phase1\_validated.pt}.
\end{center}
Its SHA-256 checksum is
\begin{center}
  \placeholder{insert SHA-256 from reports/checkpoint\_sha256.txt}.
\end{center}

\section{Principal verification commands}

\begin{lstlisting}[language=bash]
python -m tests.verify_model_interface
python -m tests.verify_reference
python -m tests.verify_neohooke_response \
  --ckpt checkpoints/phase1_validated.pt
python -m tests.verify_tangent

python - <<'PY'
from tests.verify_element import main
main("checkpoints/phase1_validated.pt")
PY

python -m tests.verify_patch \
  --ckpt checkpoints/phase1_validated.pt
python -m tests.verify_fe_comparison \
  --ckpt checkpoints/phase1_validated.pt
\end{lstlisting}

\section{Source modules}

\begin{longtable}{p{0.34\textwidth}p{0.58\textwidth}}
  \toprule
  Module & Role\\
  \midrule
  \endhead
  \texttt{src/picnn.py}
    & pICNN architecture, invariant features, positive weights and reference
      normalization.\\
  \texttt{src/neohooke.py}
    & Analytical Phase~1 material and data generation.\\
  \texttt{src/constitutive.py}
    & Autograd stress and full material tangent.\\
  \texttt{src/train.py}
    & Energy--stress training, checkpointing and refinement.\\
  \texttt{src/ngsolve\_umat.py}
    & Integration-point transfer and nonlinear FE solve.\\
  \texttt{tests/verify\_*.py}
    & V1--V6 verification programs.\\
  \bottomrule
\end{longtable}

\end{document}
